\documentclass[11pt]{article}
\usepackage{amsmath,amssymb,amsthm}
\usepackage[margin=1in]{geometry}

\newtheorem{proposition}{Proposition}

\begin{document}
This is the halving describing the Mumford representation of the order 4 element in the (2,2,4,4) family. Written entirely by GPT 5.4 Thinking (see the companion scripts halving.m and halving\_all.m).


\begin{proposition}
Let
\[
A=a^2,\qquad B=b^2,\qquad C=c^2,\qquad D=d^2,
\]
and suppose that
\[
u^2=(A-C)(A-D),\qquad
v^2=(B-C)(B-D),\qquad
w^2=(A-C)(B-C),\qquad
t^2=(A-D)(B-D),
\]
so that $uv=wt$. Let
\[
C:\ y^2=x(x+A)(x+B)(x+C)(x+D),
\]
and let
\[
P_a=(-A,0),\qquad P_b=(-B,0).
\]
Then the divisor class
\[
D:=P_a+P_b-2\infty\in J(\mathbf Q)=\operatorname{Jac}(C)(\mathbf Q)
\]
is divisible by $2$.

Define
\[
\rho=\frac{a}{b},\qquad
\sigma=\frac{A-C}{w}=\frac{w}{B-C},\qquad
\tau=\frac{A-D}{t}=\frac{t}{B-D},
\]
and let $s_1,s_2,s_3,s_4$ be the elementary symmetric polynomials in
\[
1,\ \rho,\ \sigma,\ \tau.
\]
Thus
\[
s_1=1+\rho+\sigma+\tau,
\]
\[
s_2=\rho+\sigma+\tau+\rho\sigma+\rho\tau+\sigma\tau,
\]
\[
s_3=\rho\sigma+\rho\tau+\sigma\tau+\rho\sigma\tau,
\]
\[
s_4=\rho\sigma\tau.
\]
Set
\[
\delta:=1+s_2+s_4.
\]

Then one divisor class $D'\in J(\mathbf Q)$ satisfying $2D'=D$ has Mumford
representation
\[
D'=[U(x),V(x)],
\]
where
\[
U(x)=\frac{(x+A)^2+s_2(x+A)(x+B)+s_4(x+B)^2}{\delta},
\]
i.e.
\[
U(x)=x^2+\frac{2A+s_2(A+B)+2Bs_4}{\delta}\,x
+\frac{A^2+ABs_2+B^2s_4}{\delta},
\]
and
\[
V(x)=-\frac{bv}{\delta^2}\,(\Lambda x+M),
\]
with
\[
\Lambda
=
s_1s_2-s_1s_4^2+3s_1s_4+s_2s_3s_4+3s_3s_4-s_3,
\]
and
\[
M
=
As_1s_2+2As_1s_4+As_3s_4-As_3
-Bs_1s_4^2+Bs_1s_4+Bs_2s_3s_4+2Bs_3s_4.
\]

Hence
\[
2[U(x),V(x)]=P_a+P_b-2\infty.
\]

Different choices of the square roots $\rho,\sigma,\tau$ give the other halves,
differing from $D'$ by elements of $J[2]$.
\end{proposition}

\begin{proof}
Consider the birational change of variables
\[
z=-\frac{x+A}{x+B},
\qquad
Y=\frac{(B-A)^2}{bv}\,\frac{y}{(x+B)^3}.
\]
A direct calculation shows that $C$ becomes
\[
C':\quad
Y^2=z(z+1)(z+\alpha)(z+\beta)(z+\gamma),
\]
where
\[
\alpha=\frac{A}{B}=\rho^2,\qquad
\beta=\frac{A-C}{B-C}=\sigma^2,\qquad
\gamma=\frac{A-D}{B-D}=\tau^2.
\]

Under this map,
\[
P_a\longmapsto (0,0),\qquad
P_b\longmapsto \infty',\qquad
\infty\longmapsto (-1,0).
\]
Therefore
\[
P_a+P_b-2\infty
\longmapsto
(0,0)+\infty'-2(-1,0)
\sim
(0,0)-\infty',
\]
since $(-1,0)-\infty'$ is a $2$-torsion point on $\operatorname{Jac}(C')$.

Now apply Lemma~3.1(b) to the curve $C'$, whose four nonzero branch points are
\[
1,\ \rho^2,\ \sigma^2,\ \tau^2.
\]
It follows that a half of $(0,0)-\infty'$ is represented on $C'$ by
\[
U'(z)=z^2-s_2z+s_4,
\qquad
V'(z)=(s_1s_2-s_3)z-s_1s_4.
\]

Pulling this divisor back to $C$, we substitute
\[
z=-\frac{x+A}{x+B}.
\]
Multiplying $U'(z)$ by $(x+B)^2$ and normalizing to make it monic gives
\[
U(x)=\frac{(x+A)^2+s_2(x+A)(x+B)+s_4(x+B)^2}{\delta},
\qquad
\delta=U'(-1)=1+s_2+s_4.
\]

For the $v$-polynomial, the pullback of $V'(z)$ is
\[
\widetilde V(x)
=
\frac{bv}{(B-A)^2}(x+B)^2
\Bigl(-(s_1s_2-s_3)(x+A)-s_1s_4(x+B)\Bigr).
\]
Reducing $\widetilde V(x)$ modulo $U(x)$ gives the unique linear polynomial
$V(x)$ with
\[
V(x)\equiv \widetilde V(x)\pmod{U(x)}.
\]
A straightforward reduction yields
\[
V(x)=-\frac{bv}{\delta^2}\,(\Lambda x+M),
\]
where
\[
\Lambda
=
s_1s_2-s_1s_4^2+3s_1s_4+s_2s_3s_4+3s_3s_4-s_3,
\]
and
\[
M
=
As_1s_2+2As_1s_4+As_3s_4-As_3
-Bs_1s_4^2+Bs_1s_4+Bs_2s_3s_4+2Bs_3s_4.
\]

Thus the pullback divisor $[U(x),V(x)]$ satisfies
\[
2[U(x),V(x)]=P_a+P_b-2\infty,
\]
as required.
\end{proof}

\end{document}
